\section*{Innovative Claims}
\noindent This disclosure presents [ACRONYM], a novel error correction scheme that achieves superior correction capability and lower metadata overhead than BCH codes at high bit error rates, without requiring Galois field arithmetic. The scheme combines a keyed Feistel permutation, a two-dimensional hash constraint grid, and a constraint-guided combinatorial search decoder to form a practical, tunable, and hardware-friendly ECC solution. We make the following innovative claims compared to the state of the art.

\begin{enumerate}

    \item \textbf{Hash-Constrained Two-Dimensional Grid Encoding:} [ACRONYM] introduces a novel ECC encoding structure in which data bits are arranged into a two-dimensional $N \times N$ grid and independent hash digests are computed over both the rows and columns of that grid. The row-column overlap structure creates a complete bipartite constraint graph where each grid cell is covered by exactly one row constraint and one column constraint. This dual-axis coverage means a single bit-flip corrupts exactly two hash nodes (one row, one column), and the pair of mismatched nodes acts as a direct pointer to the error location---a property that no linear algebraic code exploits.

    \item \textbf{Feistel-Based Error Diffusion for Native Burst Resilience:} Prior to grid arrangement, [ACRONYM] applies a keyed multi-round Feistel permutation with a cryptographic hash (SHA-256) as the round function. This permutation scatters any contiguous burst of $b$ consecutive bit-flips uniformly across $O(b)$ distinct rows and columns of the grid. As a result, burst errors and random errors of the same total count produce statistically identical mismatch patterns and identical correction success rates---without interleaving buffers, separate burst-correcting codes, or any additional overhead. This is a fundamental structural advantage over BCH codes, which assume independent random errors and collapse in the presence of uncorrected bursts.

    \item \textbf{Cycle-Walking Feistel for Arbitrary Block Sizes:} The Feistel permutation in [ACRONYM] uses the cycle-walking technique to produce a provably bijective permutation on integer domains of any size---not just powers of two. This ensures that [ACRONYM] applies without modification to any block length $L$, including the non-power-of-two block sizes common in real storage systems (e.g., 3,072-byte flash pages, 576-byte DRAM bursts). No prior ECC permutation/interleaving scheme provides this combination of cryptographic key-based scrambling, exact-domain bijectivity, and efficient inversion.

    \item \textbf{Matched-Neighbor Pinning Pruning:} The central algorithmic innovation of [ACRONYM]'s decoder is the matched-neighbor pinning heuristic. When searching for bit-flip combinations to correct a mismatched hash node, the decoder excludes from its candidate set any bits already covered by a \emph{matched} (correctly-verified) neighboring node in the constraint graph. These bits are ``pinned''---flipping them would corrupt an already-correct constraint. In the worst case this reduces the candidate set from $N$ bits to $O(\text{mismatched neighbors})$ bits, shrinking the combinatorial search from $\binom{N}{k}$ to $\binom{|\text{candidates}|}{k}$ at flip level $k$. Empirically, this pruning reduces total hash evaluations by approximately $47\times$ compared to BCH-equivalent algebraic decoding at 5\% BER.

    \item \textbf{Greedy Flip-Level Climbing Decoder:} The [ACRONYM] decoder implements a greedy, iterative flip-level climbing strategy. Mismatched nodes are processed in ascending order of covered-bit count (fewest bits first), with secondary ordering by number of mismatched neighbors. For each node, the decoder enumerates flip combinations of size $k = 1, 2, 3, \ldots$, scoring each by total matched nodes across the entire constraint graph, and commits the highest-scoring combination. When no combination at the current flip level improves the score, the flip level ceiling is incremented. This approach avoids the exponential worst-case complexity of global search while achieving $\geq$95\% full correction at 5\% BER on 4,096-bit blocks.

    \item \textbf{$O(1/\sqrt{L})$ Metadata Overhead Scaling:} The metadata overhead of [ACRONYM] is $2Nh$ bits, where $N = \lceil\sqrt{L}\rceil$ and $h$ is the hash width. The overhead ratio relative to the data block is $\approx 2h/\sqrt{L}$---an $O(1/\sqrt{L})$ decrease as block size grows. This is in sharp contrast to BCH codes, whose overhead ratio for a fixed bit error rate $p$ is $\Theta(p \cdot \log L)$, remaining roughly constant or growing with block size. At $L = 4{,}096$ bits with CRC-32 ($h = 32$), [ACRONYM] achieves 100\% overhead; at $L = 16{,}384$ bits with the same configuration, overhead drops to $\sim$50\%. Larger blocks naturally benefit from [ACRONYM]'s architecture.

    \item \textbf{Runtime-Tunable Overhead-Correction Tradeoff:} The overhead and correction capability of [ACRONYM] can be adjusted at deployment time by selecting the hash digest width $h$ (CRC-8: 25\% overhead; CRC-16: 50\%; CRC-32: 100\%; CRC-64: 200\%) and the row/column group size $s$ (grouping $s$ rows into one hash node reduces node count by $s$ and halves overhead per factor of $s$). These are \emph{runtime} parameters requiring no redesign of the code structure---a flexibility that BCH codes fundamentally lack, since changing $t$ requires generating an entirely new code.

    \item \textbf{No Galois Field Arithmetic:} [ACRONYM] encoding and decoding use only standard cyclic redundancy check (CRC) computations---a ubiquitous primitive present in virtually every digital system---and bitwise XOR operations. No finite-field multiplications, polynomial root-finding, or Galois field lookup tables are required. This eliminates the need for dedicated GF arithmetic hardware, reduces silicon area for memory controller ECC engines, and makes [ACRONYM] implementable in software on any general-purpose processor with CRC hardware acceleration.

    \item \textbf{Approximate Error Correction for AI Weight Memory:} Deep neural networks exhibit graceful degradation under small parameter perturbations---model accuracy drops slowly with small residual errors after approximate correction. [ACRONYM] exploits this property by providing a high-probability ($\geq$95\%) correction guarantee rather than an algebraic certainty, achieving substantially lower overhead than BCH codes at equivalent BER. For DNN weight memories storing 16-bit (FP16) or 8-bit (INT8) parameters, [ACRONYM]'s 100\% overhead at 4,096-bit blocks corresponds to one additional metadata SRAM bank per data bank, a practical cost that is lower than BCH's 174\% requirement.

    \item \textbf{Intersection-Priority Candidate Ordering:} Within the matched-neighbor pinning framework, [ACRONYM] further prioritizes ``intersection indices''---bits shared between the currently selected mismatched node and at least one other mismatched neighbor---before exploring ``free'' indices (bits not covered by any matched neighbor). This priority ordering reflects the statistical likelihood that bits implicated by multiple simultaneous constraint violations are the true error locations, guiding the search toward the correct solution with fewer evaluations.

    \item \textbf{C++/Hardware-Accelerated Implementation Path:} [ACRONYM] is designed for efficient hardware acceleration. The encoding pass (one Feistel permutation + $2N$ independent CRC computations over disjoint data ranges) is fully parallelizable: all row CRCs and all column CRCs can be computed simultaneously in a single data-streaming pass. A hardware implementation comprising $N$ parallel CRC units for rows and $N$ for columns computes all metadata in $O(N)$ clock cycles. The decoding search, while iterative, operates on small candidate sets per node (due to pinning) and requires only CRC recomputation over $O(k)$ affected nodes per flip combination. A reference C++ implementation using the pybind11 framework demonstrates the practicality of deployment.

\end{enumerate}

\section*{Intended Benefits of [ACRONYM]}

\noindent The intended benefits of [ACRONYM] are summarized below.

\begin{enumerate}

    \item \textbf{Lower ECC Overhead at High Bit Error Rates:} [ACRONYM] corrects 200 or more bit-flips in a 4,096-bit block at 100\% overhead, compared to $\sim$174\% for BCH codes at equivalent correction capability. This translates directly to reduced metadata storage area in memory controllers and reduced bandwidth for ECC metadata transmission.

    \item \textbf{Native Burst Error Resilience:} The Feistel permutation converts burst errors into random-equivalent errors with no additional overhead or latency, eliminating the need for separate interleaving hardware and simplifying system design for flash and DRAM applications where burst errors dominate.

    \item \textbf{Improved Scalability to Larger Block Sizes:} As flash page sizes and memory block sizes increase, [ACRONYM]'s overhead decreases automatically due to its $O(1/\sqrt{L})$ scaling. Larger blocks cost proportionally less to protect, unlike BCH where overhead remains roughly constant per bit regardless of block size.

    \item \textbf{Hardware Simplicity:} Replacing Galois field arithmetic with standard CRC units reduces memory controller ECC engine area and power, and enables [ACRONYM] deployment on processors and controllers that lack dedicated finite-field hardware.

    \item \textbf{Runtime Flexibility:} System designers can tune the overhead-correction tradeoff at deployment time (or even at runtime) by adjusting hash width and group size, enabling [ACRONYM] to adapt to channel conditions, power budgets, or varying BER across device lifetime.

    \item \textbf{AI Accelerator Reliability:} By providing practical, low-overhead error correction for high-density DNN weight memories, [ACRONYM] enables reliable inference at aggressive supply voltages and high storage densities, supporting the deployment of AI systems in energy-constrained and radiation-exposed environments.

    \item \textbf{Broad Applicability:} [ACRONYM] applies to any data block regardless of content, making it suitable as a drop-in ECC replacement for DRAM, NAND flash, HDD, optical storage, network transmission, and in-memory computing systems.

\end{enumerate}
